\par\indent\SourceLeaf{221}%
Uebrigens lässt sich noch zeigen, dass, wenn obige vier Grössen resp. mit \(r,s,S,R\) bezeichnet werden, allemal
\[
\frac{r\BareRoot{(ee+gg+hh)}}{s}=\pm n,\qquad\frac{R\BareRoot{(EE+GG+HH)}}{S}=\pm N
\]
\noindent wird, \(n\) und \(N\) in der Bedeutung des 5. Art. genommen; wir übergehen jedoch hier den nicht schwer zu findenden Beweis dieses Theorems, da dieses für unsern Zweck nicht weiter nöthig ist.

\par\bigskip\Needspace{7\baselineskip}
\pdfbookmark[1]{Randbemerkungen in Gauss Handexemplar}{abbildung-randbemerkungen}
\begin{center}[\textit{Randbemerkungen in} \textsc{Gauss} \textit{Handexemplar}:]\end{center}
\begingroup\small
\par\indent[Art. 10 neben der letzten Gleichung zur Bestimmung von \(\lambda\)] oder \(\lambda=\cos u^*\), wenn für \(u=u^*\) der Minimalwerth [des Vergrösserungsverhältnisses] Statt finden soll.

[Art. 12 neben der Gleichung, durch welche hier im Abdruck die Grösse \(w\) eingeführt wird] Das Zeichen \(\omega\) ist gegen meine Absicht im Druck gebraucht: es sollte \(w\) sein.

[Art. 13 neben den Gleichungen, die sich auf die durch die Function \(f\upsilon=\upsilon-i\log k\) bestimmte Abbildung beziehen, sind die entsprechenden Gleichungen für die Function \(f\upsilon=a\upsilon-i\log k\) verzeichnet, welche später in der ersten Abhandlung der Untersuchungen über Gegenstände der höhern Geodäsie aufgenommen wurden.]
\par\endgroup
\begin{center}\rule{1.2in}{.3pt}\end{center}
